\exercisesection{§ 2 的习题}

\begin{enumerate}
\def\labelenumi{\arabic{enumi})}
\tightlist
\item
  设 K 是一个有单位元的交换环，设 E 是一个 A-模，并设 \(E^{*}\) 是它的对偶。用 \(\varphi\) 表示从 \(E \otimes E^{*}\) 到 End(E) 的典范同态。
\end{enumerate}

\begin{enumerate}
\def\labelenumi{\alph{enumi})}
\tightlist
\item
  \(\operatorname{End}(\mathbf{E})\) 中任何形如下式且异于 1 的元素
\end{enumerate}

\[
s _ {x, y ^ {*}} = 1 - \varphi (x \otimes y ^ {*}),
\]

其中 \(x \in E\)，\(y^{*} \in E^{*}\)，称为 E 中的一个伪反射。这样的元素 s 称为反射，如果可以选取 x 与 \(y^{*}\) 使得 \(\langle x, y^{*} \rangle = 2\)；证明此时有 \(s^{2} = 1\) 且 \(s(x) = -x\)。

\begin{enumerate}
\def\labelenumi{\alph{enumi})}
\setcounter{enumi}{1}
\tightlist
\item
  设 \(x \in \mathbf{E}\)，\(y^{*} \in \mathbf{E}^{*}\) 满足 \(\langle x, y^{*} \rangle = 1\)。证明 \(\mathbf{E}\) 是子模 \(Kx\)（由 \(x\) 生成）与正交补 \(H\)（\(y^{*}\) 的正交补）的直和。证明 \(Kx\) 是以 \(x\) 为基的自由模，且 \(s = s_{2x,y^*}\) 在 \(H\) 上等于 1，在 \(Kx\) 上等于 -1。
\end{enumerate}

\begin{enumerate}
\def\labelenumi{\arabic{enumi})}
\setcounter{enumi}{1}
\tightlist
\item
  沿用 Exerc. 1 的记号，证明 \(\det(s_{x,y^*}) = 1 - \langle x, y^* \rangle\)（当 \(E\) 是一个有限型的自由 \(K\)-模时）。
\end{enumerate}

¶ 3) 设 V 是一个以 \(e_1, \ldots, e_l\) 为基的复 Hilbert 空间。对 \(1 \leqslant i \leqslant l\)，设 \(s_i\) 是一个酉伪反射，其向量为 \(e_i\)，使得 \(s_i(e_i) = c_i e_i\)，其中 \(c_i \neq 1\)；V 的一个元素在 \(s_i\) 下不变当且仅当它与 \(e_i\) 正交。设 W 是 \(\mathbf{GL}(V)\) 的一个子群，由诸 \(s_i\) 生成。

\begin{enumerate}
\def\labelenumi{\alph{enumi})}
\item
  设 \(i\) 是一个 \(\geqslant 1\) 的整数。证明 \(\bigwedge^{i}\mathrm{V}\) 中在 W 下不变的元素都是零。（对 \(l\) 用归纳法论证；若 \(\mathrm{V}'\) 是 V 中由 \(e_1, \ldots, e_{l-1}\) 生成的子空间，且 \(e\) 是一个与 \(\mathrm{V}'\) 正交的非零向量，则 \(\bigwedge^{i}\mathrm{V}\) 的任何元素都可以写成 \(a + (b \wedge e)\) 的形式，其中 \(a \in \bigwedge^{i}\mathrm{V}'\)，\(b \in \bigwedge^{i-1}\mathrm{V}'\)；若 \(a + (b \wedge e)\) 在 W 下不变，则 \(a\) 与 \(b\) 在子群 \(\mathrm{W}'\)（由 \(s_1, \ldots, s_{l-1}\) 生成）下不变，于是归纳假设适用。）
\item
  假设 W 有限。利用 a) 证明：对 V 的任何自同态 A，
\end{enumerate}

\[
\begin{array}{c} \sum_ {w \in \mathrm{W}} \det (\mathrm{A} - w) = \operatorname{Card} (\mathrm{W}). \det (\mathrm{A}) \\ \sum_ {w \in \mathrm{W}} \det (1 - \mathrm{A} w) = \operatorname{Card} (\mathrm{W}). \end{array}
\]

由此推出，对任何 \(A \in \text{End}(V)\)，存在 \(w \in W\) 使得 Aw 没有非零的不动点。

\begin{enumerate}
\def\labelenumi{\alph{enumi})}
\setcounter{enumi}{2}
\item
  设 \(\Gamma\) 是以 \([1,l]\) 为顶点集的图，其棱是那些集合 \(\{i,j\}\)，使得 \(e_{i}\) 与 \(e_{j}\) 不正交。证明 V 是单 W-模当且仅当 \(\Gamma\) 连通且非空。
\item
  假设 V 是一个单 W-模。证明 W-模 \(\bigwedge^{i}V\)（\(0 \leqslant i \leqslant l\)）都是单模。（证明存在一个整数 j 使得图 \(\Gamma - \{j\}\) 连通。对 l 用归纳法论证，并对子空间 \(V'\)（由诸 \(e_{i}, i \neq j\) 生成）应用归纳假设。）证明这些模两两不同构 \(^{6}\)。
\end{enumerate}

